\section{テスト理論と因子分析}
\subsection{授業内容}
\subsubsection{科目の中でのこのコマの位置づけ}
目に見えないものを測定するという意味で，テスト理論は心理学の測定と関係が深い。前回は社会心理学における態度の測定を前提に議論したが，測定に関してはテスト理論で一般的に議論できる。

真のスコアと誤差とに分解すること，誤差の基本的な仮定を確認した上で，古典的テスト理論を項目と被験者の特性に分割することで因子分析モデルに展開されるところを見る。また，多因子モデルに拡張した上で，その数理的展開から，信頼性と妥当性に言及できることを確認する。数式の展開は代数の基本的な特徴を確認すれば問題なくフォローできるはずである。

\subsubsection{コマ主題細目}
\begin{description}
	\item[古典的テスト理論]古典的テスト理論についての復習である。その基本モデルについて触れ，その平均値と分散が意味するところから測定モデルの意味するところ（誤差が相殺しあうこと）と信頼性の定義が導出できることを改めて確認しておく。

	\item[因子分析モデル]因子分析モデルは，古典的テスト理論のモデルを拡張したものである。まずは単因子モデルを例に，項目特性と被験者特性が分離されたことを確認する。その上で，性格検査や知能検査などの歴史に触れながら，多因子モデルについて解説する。多因子モデルを例に記号や添字を確認しておく。
	      \begin{flushright}
		      $\to$\textcite{kosugi2018}のPp.173--177
	      \end{flushright}

	\item[因子分析の第2定理]因子分析モデルを展開することで，相関係数が因子負荷量の積和で表現できること，因子得点がその仮定から計算上消えることを確認する。得られた指揮は因子分析の第二定理と呼ばれ，妥当性に関する議論がここから導かれることをみる。

	\item[因子分析の第1定理]ある項目自身の相関係数を考えることで，因子分析の第一定理にたどり着く。ここで共通因子の二乗和を共通性と呼ぶことにすると，信頼性の考え方が項目レベルで行われるように発展したことが確認できる。
	      \begin{flushright}
		      $\to$\textcite{kosugi2018}のPp.173--177
	      \end{flushright}
\end{description}
\subsubsection{キーワード}
\begin{itemize}
	\item 古典的テスト理論
	\item 因子分析法
	\item 因子分析の定理
\end{itemize}
\subsection{授業情報}
\paragraph{コマの展開方法}
講義
\subsubsection{予習・復習課題}
\paragraph{予習}一年時に信頼性・妥当性について，あるいは古典的テスト理論について学んだことを復習し，どのような概念であったかを再確認しておくことが望ましい。
\paragraph{復習}テスト理論と因子分析モデルの関係について，因子分析モデルはどこが新しく何を改定しようとしたのかについて，自分なりの言葉で説明できるようになろう。
